\chapter{Tendinopathy (Tendonitis)}

\section*{Definition}
Tendinopathy is a broad term encompassing tendon injury and failed healing, characterized by pain, swelling, and impaired performance. It includes tendonitis (acute inflammation) and tendinosis (chronic degeneration without significant inflammation). Common sites include the rotator cuff, lateral elbow (tennis elbow), patellar tendon, and Achilles tendon.

\section*{What Happens in Your Body}
Tendons transfer force from muscle to bone and are composed primarily of type I collagen organized in parallel fibrils. Repetitive overload beyond the tendon's mechanical capacity causes cumulative microtrauma. In acute tendonitis, there is an inflammatory infiltrate with swelling and hyperemia. Chronic tendinosis involves collagen disorganization, mucoid degeneration, neovascularization, and absence of inflammatory cells. Failed healing results from an imbalance between matrix degradation (MMP upregulation) and synthesis.

\section*{Causes}
\begin{itemize}
  \item \textbf{Overuse:} Repetitive activity in sports (running, jumping, throwing, tennis, golf) or occupation
  \item \textbf{Biomechanical:} Muscle imbalance, poor technique, improper equipment, training errors
  \item \textbf{Age-related:} Tendon vascularity and collagen turnover decrease with age
  \item \textbf{Systemic:} Diabetes mellitus, obesity, rheumatoid arthritis, gout, hypercholesterolemia
  \item \textbf{Medication-related:} Fluoroquinolone antibiotics, statins, corticosteroids (tendon rupture risk)
  \item \textbf{Anatomic:} Haglund deformity (Achilles), hypovascular zones (supraspinatus, Achilles midsubstance)
\end{itemize}

\section*{When to See a Doctor}
See your doctor if:
\begin{itemize}
  \item Pain persistent beyond 2 weeks of conservative management
  \item Tendon swelling with palpable defect or bruising (suspect rupture)
  \item Loss of function or inability to bear weight
  \item Acute onset of severe pain after a popping sensation (acute rupture)
\end{itemize}

\section*{Related Symptoms}
\begin{itemize}
  \item Localized pain over the tendon, worse with eccentric loading or stretching
  \item Palpable crepitus or thickening of the tendon
  \item Morning stiffness or start-up pain that improves with warm-up
  \item Swelling, warmth, and redness in acute tendonitis
  \item Weakness of the associated muscle group
\end{itemize}
\textbf{Important:} Tendon rupture, particularly of the Achilles, biceps, or rotator cuff, presents with a sudden pop, weakness, and a palpable gap; surgical evaluation is urgent.

\section*{Self-Care / Home Management}
\begin{itemize}
  \item Relative rest and activity modification to avoid painful loads
  \item Ice massage over the affected tendon for 10--15 minutes post-activity
  \item Eccentric strengthening exercises (e.g., Alfredson protocol for Achilles tendinopathy)
  \item Over-the-counter NSAIDs for acute pain (limited role in chronic tendinosis)
  \item Supportive bracing, taping, or heel lift as appropriate to the site
\end{itemize}

\section*{How Your Doctor Will Evaluate You}
\begin{itemize}
  \item Clinical diagnosis based on history and palpation of the symptomatic tendon
  \item Ultrasound to assess tendon thickness, fiber disruption, neovascularity, and Doppler activity
  \item MRI for detailed evaluation of tendon integrity, partial tears, and peritendinous pathology
  \item Radiography to exclude calcific tendonitis, enthesophytes, or avulsion fractures
  \item Musculoskeletal ultrasound elastography for tissue stiffness assessment (evolving)
\end{itemize}

\section*{Treatment Options}
\begin{itemize}
  \item Eccentric exercise program as the cornerstone of tendinopathy rehabilitation
  \item NSAIDs for acute tendonitis (2--4 weeks)
  \item Corticosteroid injection for short-term relief in acute tendonitis (avoid in chronic tendinosis; increases rupture risk)
  \item Platelet-rich plasma (PRP) injection for chronic tendinopathy (moderate evidence)
  \item Extracorporeal shock wave therapy for chronic calcific and non-calcific tendinopathy
  \item Percutaneous needle tenotomy or dry needling for recalcitrant cases
  \item Surgical removal of dead tissue and repair for complete rupture or failed non-operative management
\end{itemize}

\clearpage
